\documentclass{article}
\usepackage[UTF8]{ctex}
\usepackage{amsmath,amssymb,amsfonts}
\usepackage{enumitem}
\usepackage{geometry}
\geometry{a4paper, margin=1in}
% === 错题本排版优化（不参与题目标识） ===
\setlength{\parindent}{0pt}
\setlength{\parskip}{0.35em plus 0.10em minus 0.08em}
\setlength{\emergencystretch}{3em}
\setlist{itemsep=0.25em, topsep=0.40em, parsep=0pt}
\allowdisplaybreaks[2]
% === 错题本排版优化结束 ===
\begin{document}
\section*{首都师范大学2025年数学分析真题}
\section*{32．首都师范大学2025年数学分析试卷}
一．（24 分）计算题．

\begin{itemize}

\subsection{题目1}
  \item[1．] （8 分）求极限 \(\lim _{x \rightarrow 0}\left(\frac{\sin x}{x}\right)^{\frac{1}{1-\cos x}}\) ．

\subsection{题目2}
  \item[2．] （8 分）设函数 \(f(x)\) 在点 \(a\) 可导且 \(f(a) \neq 0\) ，求 \(\lim _{n \rightarrow \infty}\left[\frac{f\left(a+\frac{1}{n}\right)}{f(a)}\right]^{n}\) ．

\subsection{题目3}
  \item[3．] （8 分）求级数 \(\sum_{n=1}^{\infty} \frac{n}{3^{n}}\) 的和。
\end{itemize}

二．（15 分）考虑方程

\[
e^{x+2 y+3 z}+x y z=1
\]

\begin{itemize}

  \item[1．] （7 分）设函数 \(z=z(x, y)\) 由（？？）确定，求 \(\mathrm{d} z\) ．

  \item[2．] （8 分）设曲面 \(S\) 由（？？）确定，求 \(S\) 上过点 \((0,0,0)\) 的切平面方程和法线方程．
\end{itemize}

三．（16 分）考虑函数 \(f(x)= \begin{cases}\frac{x}{3}+x^{2} \sin \frac{1}{x}, & x \neq 0 ; \\ 0, & x=0 .\end{cases}\)\\

1．（8 分）讨论函数 \(f(x)\) 的连续性和可导性。\\

2．（8 分）讨论是否存在 \(x=0\) 的一个邻域，使得 \(f(x)\) 在该邻域内单调并给出理由．

四．（10 分）证明不等式 \(x y \leq x \ln x-x+e^{y}, x \geq 1, y \geq 0\) ．

五。（12 分）设 \(L\) 为球面 \(x^{2}+y^{2}+z^{2}=a^{2}(a>0)\) 和平面 \(x+y+z=0\) 的交线，一质点在点 \(M(x, y, z)\)处受到力 \(\mathbf{F}=y \mathbf{i}+z \mathbf{j}+x \mathbf{k}\) 的作用，其中 \(\mathbf{i}=(1,0,0), \mathbf{j}=(0,1,0), \mathbf{k}=(0,0,1)\) ，此质点沿 \(L\) 运动一周，方向从 \(z\) 轴正项看为逆时针方向，求 \(\mathbf{F}\) 所作的功 \(W\) 。

六．（10分）设 \(f(x)\) 是区间 \([a, b]\) 上的连续函数，且满足条件 \(\int_{a}^{b} x^{k} f(x) \mathrm{d} x=0, k=0,1, \cdots, n\) ．证明：函数 \(f(x)\) 在 \((a, b)\) 内至少有 \(n+1\) 个不同的零点．

\begin{itemize}
  \item[七．] （20 分）设 \(f(x)=\sum_{n=1}^{\infty} \frac{\cos n x}{n^{2}+1}\) ．

  \item[1．] （10 分）证明：\(f(x)\) 在区间 \((0,2 \pi)\) 内连续，且有连续的导函数．

  \item[2．] （10 分）记 \(F(x)=\int_{0}^{x} f(t) \mathrm{d} t\) ，证明：\(-\frac{1}{30}<F\left(\frac{\pi}{2}\right)<\frac{1}{2}\) ．
\end{itemize}

\subsection{题目902}
原题边界“八．”：（18分）讨论下列反常积分的敛散性（包括绝对收敛，条件收敛和发散）
\[
\int_{0}^{+\infty} \frac{x \sin x}{1+x^{p}} \mathrm{~d} x, p \geq 0
\]

九．（15 分）设 \(f(x)\) 以 \(2 \pi\) 为周期且具有一阶连续的导函数．证明：\(f(x)\) 的傅里叶级数在 \((-\infty,+\infty)\)上一致收敛于 \(f(x)\) 。

十．（10 分）若函数 \(f(x)\) 在 \((0,1)\) 内可微，且 \(\lim _{x \rightarrow 0^{+}} f(x)=\infty\) ，试证： \(\lim _{x \rightarrow 0^{+}} f^{\prime}(x)=\infty\) ．
\end{document}
